\chapter{常微分方程}
\section{一阶方程：先识别，再变形}
\ruleitem{可分离变量}{$y'=p(x)q(y)$：当 $q(y)\ne0$，积分 $\int\dd y/q(y)=\int p(x)\dd x+C$。每个满足 $q(y_0)=0$ 的常值函数须单独检查，除法可能丢解。}
\ruleitem{齐次型}{$y'=F(y/x)$，在 $x\ne0$ 的区间令 $y=vx$，则 $xv'=F(v)-v$，再分离。这里“齐次型”与线性齐次方程是不同概念。}
\ruleitem{一阶线性}{$y'+P(x)y=Q(x)$，连续系数区间内
\[
\mu(x)=\exp\left(\int P(x)\dd x\right),\qquad
y=\mu^{-1}\left(\int\mu Q\dd x+C\right).
\]
公式来自 $(\mu y)'=\mu Q$；积分常数只保留一个独立常数即可。}
\ruleitem{Bernoulli}{$y'+Py=Qy^m$，$m\ne0,1$，在允许除以 $y^m$ 的区间令 $u=y^{1-m}$，得 $u'+(1-m)Pu=(1-m)Q$。检查 $y=0$ 是否为原方程解，非整数幂还需检查实数定义域。}
\ruleitem{恰当微分}{$M(x,y)\dd x+N(x,y)\dd y=0$。在单连通域内 $M,N\in C^1$ 且 $M_y=N_x$，存在势函数 $\Phi$，$\Phi_x=M,\Phi_y=N$，解为 $\Phi=C$。先对一个变量积分，再用另一个偏导确定遗漏的单变量函数。}
\ruleitem{一元积分因子}{若 $(M_y-N_x)/N$ 只依赖 $x$，可取 $\mu(x)=\ee^{\int(M_y-N_x)/N\dd x}$；若 $(N_x-M_y)/M$ 只依赖 $y$，可取相应的 $\mu(y)$。除法处仍需检查非零条件。}
\ruleitem{初值解唯一性}{对 $y'=F(x,y)$，$F$ 连续且关于 $y$ 局部 Lipschitz（例如 $F_y$ 连续），给定初值可保证局部存在唯一解。仅连续通常只保证存在，不能保证唯一。解的区间可能受系数奇点、爆破解限制。}

\section{二阶常系数、降阶与共振}
\ruleitem{线性结构}{$y''+py'+qy=r(x)$ 的通解为 $y_h+y_p$。两个非齐次解之差是齐次解；非齐次解随意相加一般不再是该方程的解。}
特征方程 $\lambda^2+p\lambda+q=0$：
\[
y_h=\begin{cases}
C_1\ee^{\lambda_1x}+C_2\ee^{\lambda_2x},&\lambda_1\ne\lambda_2\text{ 为实根},\\
(C_1+C_2x)\ee^{\lambda x},&\lambda\text{ 为二重实根},\\
\ee^{\alpha x}(C_1\cos\beta x+C_2\sin\beta x),&\alpha\pm\ii\beta,\ \beta\ne0.
\end{cases}
\]
若右端为 $\ee^{\alpha x}P_m(x)$，试取 $y_p=x^k\ee^{\alpha x}Q_m(x)$，$k$ 为 $\alpha$ 在特征方程中的重数。若右端含 $\cos\beta x,\sin\beta x$，两者都要保留，按复根 $\alpha+\ii\beta$ 的重数乘 $x^k$。所谓共振，就是试解与齐次解重合后需要补乘 $x$。
\ruleitem{常用降阶}{不显含 $y$ 的 $y''=F(x,y')$，令 $p=y'$；不显含 $x$ 的 $y''=F(y,y')$，令 $p=p(y)$，$y''=p\dd p/\dd y$，同时检查常值解。}
\ruleitem{Euler 方程}{$x^2y''+axy'+by=r(x)$，在 $x>0$ 令 $t=\ln x$，有 $xy'=Y_t$，$x^2y''=Y_{tt}-Y_t$，化为常系数方程；$x<0$ 可用 $t=\ln|x|$ 分区间处理。}

\section{结构识别真题}
\problem{43}{由三个非齐次解反求方程}{2009 年第 1 届，非数学类，五，\src{01}}
同一二阶常系数非齐次线性方程有三个解：$y_1=x\ee^x+\ee^{2x}$，$y_2=x\ee^x+\ee^{-x}$，$y_3=x\ee^x+\ee^{2x}-\ee^{-x}$。求以最高阶系数为 1 的方程及通解。
\sol $y_1-y_3=\ee^{-x}$，$y_2-y_3=2\ee^{-x}-\ee^{2x}$ 是两个线性无关齐次解，故特征根为 $-1,2$，齐次算子为 $D^2-D-2$。
取特解 $x\ee^x$，代入得 $r(x)=(1-2x)\ee^x$。所以
\[
\ans{y''-y'-2y=(1-2x)\ee^x},\qquad
\ans{y=C_1\ee^{2x}+C_2\ee^{-x}+x\ee^x}.
\]
从解的差得到齐次结构，再找右端，顺序不能颠倒。

\problem{44}{重根对应的两个基本解}{2014 年第 6 届，非数学类，一(1)，\src{06}}
已知 $\ee^x,x\ee^x$ 是某二阶常系数齐次方程的两个解，求方程。
\sol 两解线性无关，对应特征根 1 的重数为 2。因此规范化方程为 $\ans{y''-2y'+y=0}$。只用 $\ee^x$ 不能唯一确定两个特征根；第二个解 $x\ee^x$ 提供重根信息。

\problem{45}{常数右端与初值}{第 14 届初赛补赛二，一(3)，\src{14}}
解 $y''-2y'-3y=1$，$y(0)=0,y'(0)=1$。
\sol 齐次根为 $3,-1$，常数特解为 $-1/3$，故 $y=C_1\ee^{3x}+C_2\ee^{-x}-1/3$。初值给出 $C_1+C_2=1/3$、$3C_1-C_2=1$，解得 $C_1=1/3,C_2=0$，因此 $\ans{y=(\ee^{3x}-1)/3}$。代回原方程和两个初值复核。

\problem{46}{微分形式中识别 y 平方}{2024 年第 16 届，非数学 A 类二／B 类四，共题，\src{16A}\src{16B}}
求 $(x^3-y^2)\dd x+(x^2y+xy)\dd y=0$ 的通解。
\method{只有 $y^2$ 和 $y\dd y$，令 $u=y^2$ 变为一阶线性。}
\sol 以 $x$ 为自变量，$\dd u=2y\dd y$，得
\[x^3-u+\frac{x(x+1)}2u'=0,\qquad
u'-\frac2{x(x+1)}u=-\frac{2x^2}{x+1}.\]
在 $x\ne0,-1$ 的连通区间，积分因子 $\mu=((x+1)/x)^2$，所以
\[(\mu u)'=-2(x+1),\qquad
\ans{\frac{(x+1)^2}{x^2}y^2+(x+1)^2=C}.\]
也可整理为 $(x+1)^2(x^2+y^2)=Cx^2$。后式只是方便表达，并不自动把被除掉的直线恢复为图像解；若讨论微分形式的所有积分曲线，竖直直线 $x=0$、$x=-1$ 也须直接代回原式检查，它们确为积分曲线。过奇点的局部解需按原式分析。

\problem{47}{三角非线性用 sec 代换}{2025 年第 17 届，非数学 A 类，三，\src{17A}}
解 $x\ln x\sin y\,y'+(1-x\cos y)\cos y=0$（实数范围 $x>0$）。
\sol 在 $\cos y\ne0$ 的区间令 $u=\sec y$。除以 $x\cos^2y$ 得
\[\ln x\,u'+\frac ux=1,\qquad (u\ln x)'=1.\]
所以 $\ans{\sec y\ln x=x+C}$，或隐式写成 $\ln x=(x+C)\cos y$。在 $x\ne1$ 的常规区间可按一阶线性求解；写成乘积求导则避免了额外除以 $\ln x$。
另有被 $\cos y$ 除法遗漏的常值解 $\ans{y=\pi/2+k\pi}$（$k\in\mathbb Z$）。隐式式子在 $x=1$ 可能退化，实际延拓必须满足原方程，不能只因代入隐式等式成立便认定为可导解。
